\section{FHE Background}

\label{sec:background}
% ============================================================================

\begin{definition}[Fully Homomorphic Encryption]
An FHE scheme $\mathcal{E} = (\texttt{KeyGen}, \texttt{Enc}, \texttt{Dec},
\texttt{Eval})$ satisfies: for any function $f$ and plaintexts $x_1, \ldots, x_k$,
\begin{equation}
\texttt{Dec}_{sk}\!\left(\texttt{Eval}_{ek}(f, \texttt{Enc}_{pk}(x_1), \ldots,
\texttt{Enc}_{pk}(x_k))\right) = f(x_1, \ldots, x_k),
\end{equation}
where $(pk, sk, ek) \leftarrow \texttt{KeyGen}(\lambda)$ for security
parameter $\lambda$.
\end{definition}

The computational cost of FHE depends on the depth of the arithmetic circuit
being evaluated. Each homomorphic operation increases noise in the ciphertext;
\emph{bootstrapping} refreshes noise levels but is expensive.

\subsection{CKKS for Neural Networks}

CKKS~\cite{ckks2017} supports approximate arithmetic on encrypted real-valued
vectors. A single ciphertext encodes $N/2$ complex values (or $N$ real values)
via SIMD packing, where $N$ is the polynomial ring degree.

\begin{definition}[CKKS Encoding]
Given a plaintext vector $\mathbf{z} \in \mathbb{C}^{N/2}$, the encoding map
$\sigma^{-1}: \mathbb{C}^{N/2} \to \mathbb{Z}[X]/(X^N+1)$ produces a
polynomial $m(X)$ such that $\sigma(m) \approx \Delta \cdot \mathbf{z}$, where
$\Delta$ is a scaling factor controlling precision.
\end{definition}

CKKS supports:
\begin{itemize}[leftmargin=1.1em]
    \item \textbf{Addition}: $\texttt{Enc}(a) + \texttt{Enc}(b) =
    \texttt{Enc}(a + b)$, consuming no multiplicative depth.
    \item \textbf{Multiplication}: $\texttt{Enc}(a) \cdot \texttt{Enc}(b) =
    \texttt{Enc}(a \cdot b)$, consuming one level of multiplicative depth.
    \item \textbf{Rotation}: $\texttt{Rot}(\texttt{Enc}(\mathbf{z}), k) =
    \texttt{Enc}(\text{rotate}(\mathbf{z}, k))$, enabling SIMD-parallel
    operations.
\end{itemize}

A neural network layer $\mathbf{y} = \phi(W\mathbf{x} + \mathbf{b})$ is
evaluated homomorphically by encoding matrix-vector products as diagonal
rotations and approximating activation functions with low-degree polynomials.

\begin{proposition}[Polynomial Activation Approximation]
The ReLU function $\phi(x) = \max(0, x)$ is approximated on $[-B, B]$ by a
degree-$d$ minimax polynomial $\hat{\phi}_d(x)$ with error:
\begin{equation}
\|\phi - \hat{\phi}_d\|_\infty \leq \frac{2B}{\pi} \cdot
\frac{1}{2^d},
\end{equation}
where $d = 7$ suffices for $<$0.1\% classification accuracy loss.
\end{proposition}

\subsection{TFHE for Decision Trees}

TFHE~\cite{tfhe2020} operates on encrypted bits and integers with exact
arithmetic. Programmable bootstrapping allows evaluating arbitrary lookup tables
on encrypted values in a single operation.

\begin{definition}[Programmable Bootstrapping]
Given a ciphertext $c$ encrypting $m \in \mathbb{Z}_p$ and a function $f:
\mathbb{Z}_p \to \mathbb{Z}_p$, programmable bootstrapping produces a fresh
ciphertext encrypting $f(m)$ with reset noise:
\begin{equation}
\texttt{PBS}(c, f) = \texttt{Enc}(f(\texttt{Dec}(c))),
\end{equation}
without revealing $m$ to the evaluator.
\end{definition}

Decision tree evaluation reduces to a sequence of encrypted comparisons. Each
tree node performs an encrypted comparison $\texttt{Enc}(x_i) \stackrel{?}{<}
t_j$ using TFHE's Boolean gates, and the result selects the left or right child.

\begin{theorem}[Decision Tree Complexity]
A decision tree of depth $d$ with threshold values in $\mathbb{Z}_{2^b}$
requires $d$ encrypted comparisons, each costing $b$ programmable bootstrappings.
Total cost: $O(d \cdot b)$ bootstrappings.
\end{theorem}

For a gradient-boosted ensemble with $T$ trees of depth $d$, the total cost is
$O(T \cdot d \cdot b)$ bootstrappings. Crucially, all $T$ trees can be
evaluated in parallel.

% ============================================================================
